\documentclass[aspectratio=169]{beamer}
\setbeamerfont{headline}{size=\footnotesize}

\mode<presentation>{
    \usetheme{Montpellier}
    \setbeamercovered{transparent}
    \usecolortheme{beaver}
}

\usepackage[english,brazil]{babel}
\usepackage[utf8]{inputenc}
\usepackage{fontawesome}
\usepackage{graphicx}
\usepackage{listings}
\usepackage{url}
\usepackage{colortbl}
\usepackage{caption}
\usepackage{booktabs}
\usepackage{ragged2e}

\definecolor{mineblue}{RGB}{0,90,180}
\definecolor{mineorange}{RGB}{230,115,0}
\definecolor{minegreen}{RGB}{34,139,34}
\definecolor{minered}{RGB}{180,20,40}
\definecolor{minegray}{RGB}{100,100,100}

\usepackage{tikz}
\usetikzlibrary{arrows.meta,shapes.geometric,positioning,fit,backgrounds,calc}
\usepackage{pgfplots}
\pgfplotsset{compat=1.18}

\definecolor{codegreen}{rgb}{0,0.6,0}
\definecolor{codegray}{rgb}{0.5,0.5,0.5}
\definecolor{codepurple}{rgb}{0.58,0,0.82}
\definecolor{backcolour}{rgb}{0.95,0.95,0.92}

\newcommand{\reflexao}[1]{\begin{alertblock}{Reflexão}#1\end{alertblock}}
\newcommand{\key}[1]{\underline{#1}}

\title[Mineração de Dados]{Mineração de Dados: Conceitos e Aplicações}
\subtitle{Módulo 4: Técnicas de Regressão}
\author[Elton Sarmanho]{Prof. Elton Sarmanho\inst{1}\\\footnotesize \href{mailto:eltonss@ufpa.br}{eltonss@ufpa.br}}
\institute[UFPA]{\inst{1} Faculdade de Sistemas de Informação - UFPA Campus Cametá}
\date{\today}

\tikzstyle{block} = [rectangle, draw, fill=blue!20, text width=5em, text centered, rounded corners, minimum height=4em]
\tikzstyle{line} = [draw, -latex']
\tikzstyle{cloud} = [draw, ellipse, fill=red!20, node distance=3cm, minimum height=2em]

\begin{document}

% ============================================================
% 1. TÍTULO
% ============================================================
\begin{frame}
    \titlepage
    \begin{center}
        \tiny \faGithub \ \href{https://github.com/eltonsarmanho}{github.com/eltonsarmanho} \quad
        \faLinkedinSquare \ \href{https://www.linkedin.com/in/elton-sarmanho-836553185/}{Elton Sarmanho}
    \end{center}
\end{frame}

% ============================================================
% 2. ROTEIRO
% ============================================================
\begin{frame}{Roteiro da Aula}
    \tableofcontents
\end{frame}

% ============================================================
\section{O Problema de Regressão}
% ============================================================

\begin{frame}{Classificação vs.\ Regressão}
\begin{columns}
\column{0.55\textwidth}
\begin{table}
\centering\small
\begin{tabular}{>{\bfseries}p{2.0cm}p{2.0cm}p{2.0cm}}
\toprule
Aspecto & Classificação & Regressão \\
\midrule
Saída & Discreta & Contínua \\
Exemplo & Spam/Ham & Preço \\
Métrica & Acurácia, F1 & RMSE, R² \\
Perda & Cross-entropy & MSE, MAE \\
\bottomrule
\end{tabular}
\end{table}
\column{0.42\textwidth}
\begin{block}{Exemplos Práticos}
\begin{itemize}\small
  \item \textbf{Preço de imóvel}: área, localização
  \item \textbf{Temperatura}: dados climáticos
  \item \textbf{Demanda}: estoque e logística
  \item \textbf{Risco financeiro}: spreads, VaR
  \item \textbf{Dosagem médica}: peso, idade
\end{itemize}
\end{block}
\end{columns}
\vspace{0.3em}
\begin{alertblock}{}
\centering\small
A saída é \textbf{contínua}: $y \in \mathbb{R}$, não uma classe discreta.
\end{alertblock}
\end{frame}

\begin{frame}{Tipos de Regressão}
\begin{columns}
\column{0.5\textwidth}
\begin{block}{Modelos Lineares}
\begin{itemize}
  \item \textbf{Linear Simples}: $\hat{y} = \beta_0 + \beta_1 x$
  \item \textbf{Múltipla}: $\hat{y} = \boldsymbol{\beta}^T\mathbf{x}$
  \item \textbf{Polinomial}: inclui $x^2, x^3, \ldots$
\end{itemize}
\end{block}
\vspace{0.5em}
\begin{block}{Regularizada}
\begin{itemize}
  \item \textbf{Ridge}: penalidade L2
  \item \textbf{LASSO}: penalidade L1 (esparsidade)
  \item \textbf{ElasticNet}: L1 + L2
\end{itemize}
\end{block}
\column{0.47\textwidth}
\begin{exampleblock}{Não-Linear / Ensemble}
\begin{itemize}
  \item \textbf{SVR}: kernel trick
  \item \textbf{Random Forest}: média de árvores
  \item \textbf{XGBoost / LightGBM}: boosting
  \item \textbf{Redes Neurais}: MLP com saída linear
\end{itemize}
\end{exampleblock}
\vspace{0.5em}
\begin{alertblock}{\small Consulte o Material Didático}
\small Módulo 4, Seções 1 e 2
\end{alertblock}
\end{columns}
\end{frame}

% ============================================================
\section{Regressão Linear}
% ============================================================

\begin{frame}[t]{Regressão Linear Simples}
\begin{columns}
\column{0.42\textwidth}
\small
\textbf{Objetivo:} encontrar a reta que melhor ajusta os dados.

\vspace{0.5em}
$\hat{y} = \beta_0 + \beta_1 x$

\vspace{0.4em}
Estimadores OLS (Mínimos Quadrados):
\[
  \hat{\beta}_1 = \frac{S_{xy}}{S_{xx}}, \quad \hat{\beta}_0 = \bar{y} - \hat{\beta}_1\bar{x}
\]

\vspace{0.4em}
Minimiza $\displaystyle\sum_i (y_i - \hat{y}_i)^2$

\vspace{0.4em}
\begin{exampleblock}{\small Resíduos}
\small $\varepsilon_i = y_i - \hat{y}_i$ (linhas laranjas)
\end{exampleblock}

\column{0.56\textwidth}
\begin{tikzpicture}
\begin{axis}[
  width=8cm, height=5cm,
  xlabel={$x$}, ylabel={$y$},
  title={\small Regressão Linear Simples com Resíduos},
  grid=major, grid style={minegray!25},
  xmin=0, xmax=11, ymin=0, ymax=20,
  title style={font=\small}
]
\addplot[only marks, mineblue, mark size=2.5pt] coordinates {
  (1,2.1)(2,3.8)(3,5.2)(4,7.1)(5,8.9)(6,10.5)(7,12.1)(8,14.3)(9,15.7)(10,17.5)
  (1.5,3.0)(3.5,6.5)(5.5,9.8)(7.5,13.0)(9.5,16.5)
};
\addplot[thick, minered, domain=0.5:10.5, samples=50] {1.8*x + 0.3};
\addlegendentry{$\hat{y} = 1{,}8x + 0{,}3$}
\addplot[dashed, mineorange, thin] coordinates {(3,5.2)(3,5.7)};
\addplot[dashed, mineorange, thin] coordinates {(7,12.1)(7,12.9)};
\addplot[dashed, mineorange, thin] coordinates {(5,8.9)(5,9.3)};
\node[font=\tiny, mineorange] at (axis cs:3.5,5.45) {$\varepsilon_i$};
\end{axis}
\end{tikzpicture}
\end{columns}
\end{frame}

\begin{frame}[shrink=5]{Métricas de Avaliação}
\begin{table}
\centering\small
\begin{tabular}{>{\bfseries}p{2.0cm}p{4.5cm}p{4.0cm}}
\toprule
Métrica & Fórmula & Quando usar \\
\midrule
MSE  & $\frac{1}{n}\sum(y_i-\hat{y}_i)^2$ & Penaliza outliers; otimização \\
RMSE & $\sqrt{MSE}$ & Mesma unidade de $y$ \\
MAE  & $\frac{1}{n}\sum|y_i-\hat{y}_i|$ & Robusto a outliers \\
R²   & $1 - \frac{\sum(y_i-\hat{y}_i)^2}{\sum(y_i-\bar{y})^2}$ & Proporção da variância explicada \\
R² Ajust. & $1-(1-R^2)\frac{n-1}{n-p-1}$ & Penaliza features desnecessárias \\
\bottomrule
\end{tabular}
\end{table}
\vspace{0.3em}
\begin{columns}
\column{0.48\textwidth}
\begin{alertblock}{\small Atenção ao R²}
\small Sempre cresce com mais features! Use R² ajustado para comparar modelos com diferente número de variáveis.
\end{alertblock}
\column{0.48\textwidth}
\begin{exampleblock}{\small Dica Prática}
\small MAE é preferível quando há outliers. RMSE é preferível quando erros grandes são muito custosos.
\end{exampleblock}
\end{columns}
\end{frame}

% ============================================================
\section{Regularização}
% ============================================================

\begin{frame}{Por que Regularizar?}
\begin{columns}
\column{0.55\textwidth}
\begin{block}{O Problema: Overfitting}
\begin{itemize}
  \item Modelo complexo memoriza os dados de treino
  \item Erro baixo no treino, alto no teste
  \item Coeficientes muito grandes, tornando-os instáveis
\end{itemize}
\end{block}
\vspace{0.5em}
\begin{alertblock}{Solução: Penalidade}
Adicionar um termo à função de perda que \textbf{penaliza coeficientes grandes}:
\[
\text{Perda} + \lambda \cdot \text{Penalidade}(\boldsymbol{\beta})
\]
\end{alertblock}
\column{0.42\textwidth}
\begin{exampleblock}{Benefícios}
\begin{itemize}\small
  \item Generalização melhor
  \item Reduz variância
  \item LASSO: seleção de features
  \item Ridge: lida com multicolinearidade
\end{itemize}
\end{exampleblock}
\vspace{0.5em}
\begin{block}{\small $\lambda$: hiperparâmetro}
\small $\lambda \uparrow$ mais regularização\\
$\lambda = 0$: sem penalidade (OLS)\\
Escolher via validação cruzada
\end{block}
\end{columns}
\end{frame}

\begin{frame}[t]{Ridge vs.\ LASSO}
\begin{columns}
\column{0.38\textwidth}
\small
\begin{block}{Ridge (L2)}
\small $\|\mathbf{y}-X\boldsymbol{\beta}\|^2 + \lambda\|\boldsymbol{\beta}\|_2^2$\\[0.3em]
\textbf{Efeito:} encolhe todos os coeficientes, mas \textbf{nenhum é zerado exatamente}
\end{block}
\vspace{0.4em}
\begin{block}{LASSO (L1)}
\small $\|\mathbf{y}-X\boldsymbol{\beta}\|^2 + \lambda\|\boldsymbol{\beta}\|_1$\\[0.3em]
\textbf{Efeito:} produz \textbf{soluções esparsas}: alguns coeficientes $= 0$ (seleção de features)
\end{block}
\column{0.6\textwidth}
\centering
\scalebox{0.75}{%
\begin{tikzpicture}[scale=1.4]
% L2 ball
\begin{scope}[xshift=-2.8cm]
  \draw[mineblue, thick, fill=mineblue!10] (0,0) circle (1);
  \draw[->, minegray, thin] (-1.4,0) -- (1.4,0) node[right,font=\tiny]{$\beta_1$};
  \draw[->, minegray, thin] (0,-1.4) -- (0,1.4) node[above,font=\tiny]{$\beta_2$};
  \draw[minered, dashed, thin] (0.7, 0.714) ellipse (0.5 and 0.35);
  \draw[minered, dashed, thin] (0.7, 0.714) ellipse (0.85 and 0.6);
  \filldraw[minered] (0.7, 0.714) circle(0.05) node[above right,font=\tiny,minered]{ótimo};
  \node[below,font=\tiny,align=center] at (0,-1.6) {\textbf{Ridge (L2)}\\$\|\beta\|_2 \leq t$};
\end{scope}
% L1 diamond
\begin{scope}[xshift=2.0cm]
  \draw[mineorange, thick, fill=mineorange!10] (1,0)--(0,1)--(-1,0)--(0,-1)--cycle;
  \draw[->, minegray, thin] (-1.4,0) -- (1.4,0) node[right,font=\tiny]{$\beta_1$};
  \draw[->, minegray, thin] (0,-1.4) -- (0,1.4) node[above,font=\tiny]{$\beta_2$};
  \draw[minered, dashed, thin] (0.3, 1) ellipse (0.65 and 0.45);
  \draw[minered, dashed, thin] (0.3, 1) ellipse (1.1 and 0.8);
  \filldraw[minered] (0, 1) circle(0.05) node[right,font=\tiny,minered,align=left]{ótimo\\(esparsidade!)};
  \node[below,font=\tiny,align=center] at (0,-1.6) {\textbf{LASSO (L1)}\\$\|\beta\|_1 \leq t$};
\end{scope}
\end{tikzpicture}%
}

\vspace{0.2em}
\tiny O ótimo do LASSO toca o losango num \textbf{vértice} onde $\beta_1=0$, gerando esparsidade.
\end{columns}
\end{frame}

\begin{frame}{ElasticNet e Quando Usar Cada Método}
\begin{block}{ElasticNet = Ridge + LASSO}
\[
\hat{\boldsymbol{\beta}}_{EN} = \arg\min_{\boldsymbol{\beta}} \left\{ \|\mathbf{y}-X\boldsymbol{\beta}\|^2 + \lambda_1\|\boldsymbol{\beta}\|_1 + \lambda_2\|\boldsymbol{\beta}\|_2^2 \right\}
\]
\end{block}
\vspace{0.3em}
\begin{table}
\centering\small
\begin{tabular}{>{\bfseries}p{2.5cm}p{4.0cm}p{4.0cm}}
\toprule
Método & Use quando\ldots & Vantagem \\
\midrule
Ridge & Muitas features correlacionadas & Coeficientes estáveis \\
LASSO & Quer seleção automática de features & Modelo esparso / interpretável \\
ElasticNet & Grupos de features correlacionadas & Combina estabilidade + esparsidade \\
\bottomrule
\end{tabular}
\end{table}
\vspace{0.3em}
\begin{exampleblock}{\small Dica}
\small Sempre escolha $\lambda$ por validação cruzada (\texttt{RidgeCV}, \texttt{LassoCV}, \texttt{ElasticNetCV} no scikit-learn).
\end{exampleblock}
\end{frame}

% ============================================================
\section{Trade-off Viés-Variância}
% ============================================================

\begin{frame}[t]{Trade-off Viés-Variância}
\begin{columns}
\column{0.4\textwidth}
\small
\textbf{Decomposição do erro:}
\[
\text{Erro} = \text{Viés}^2 + \text{Variância} + \sigma^2
\]

\vspace{0.4em}
\begin{itemize}\small
  \item \textbf{Underfitting:} modelo simples, alto viés, baixa variância
  \item \textbf{Overfitting:} modelo complexo, baixo viés, alta variância
  \item \textbf{Ótimo:} ponto que minimiza o erro total
\end{itemize}

\vspace{0.4em}
\begin{alertblock}{\small Implicação}
\small Regularização aumenta o viés mas reduz a variância, melhorando a generalização em geral.
\end{alertblock}

\column{0.58\textwidth}
\begin{tikzpicture}
\begin{axis}[
  width=8cm, height=5cm,
  xlabel={Complexidade do Modelo},
  ylabel={Erro},
  title={\small Trade-off Viés-Variância},
  grid=major, grid style={minegray!20},
  xmin=0, xmax=10, ymin=0, ymax=6,
  xtick={0,...,10}, ytick={0,1,...,6},
  xticklabels={,,Simples,,,,,,,,Complexo},
  legend pos=north east,
  legend style={font=\tiny},
  title style={font=\small}
]
\addplot[thick, mineblue, smooth, domain=0.5:9.5] {5/(x^0.7 + 1)};
\addlegendentry{Viés$^2$}
\addplot[thick, mineorange, smooth, domain=0.5:9.5] {0.05 * x^1.8};
\addlegendentry{Variância}
\addplot[thick, minegreen, smooth, dashed, domain=0.5:9.5] {5/(x^0.7+1) + 0.05*x^1.8 + 0.5};
\addlegendentry{Erro Total}
\addplot[only marks, minered, mark=*, mark size=3pt] coordinates {(4.2, 2.3)};
\node[font=\tiny, minered, above right] at (axis cs:4.2,2.3) {Ótimo};
\addplot[gray, dashed, domain=0:10] {0.5};
\node[font=\tiny, gray, right] at (axis cs:9.5,0.5) {Ruído};
\end{axis}
\end{tikzpicture}
\end{columns}
\end{frame}

% ============================================================
\section{Regressão Polinomial}
% ============================================================

\begin{frame}{Regressão Polinomial}
\begin{columns}
\column{0.5\textwidth}
\begin{block}{Ideia}
Adicionar $x^2, x^3, \ldots$ como novas features:
\[
\phi(\mathbf{x}) = [1,\; x,\; x^2,\; \ldots,\; x^d]^T
\]
O modelo é \textbf{linear nos parâmetros}, resolvido por OLS.
\end{block}

\vspace{0.4em}
\begin{block}{Grau $d$: como escolher?}
\begin{itemize}\small
  \item Grau baixo: underfitting (alto viés)
  \item Grau alto: overfitting (alta variância)
  \item \textbf{Solução:} validação cruzada
\end{itemize}
\end{block}

\column{0.47\textwidth}
\begin{alertblock}{Risco com Grau Alto}
Polinômio de grau $n-1$ interpola \textbf{todos} os pontos do treino com erro zero, mas generaliza pessimamente!

\vspace{0.3em}
\small Oscilações de Runge: o polinômio oscila violentamente entre os pontos de treino.
\end{alertblock}

\vspace{0.4em}
\begin{exampleblock}{\small Alternativa}
\small Usar regularização Ridge/LASSO junto com termos polinomiais: controla as oscilações.
\end{exampleblock}
\end{columns}
\end{frame}

% ============================================================
\section{Métodos Ensemble para Regressão}
% ============================================================

\begin{frame}{Random Forest Regressor}
\begin{columns}
\column{0.5\textwidth}
\begin{block}{Como funciona}
\begin{itemize}\small
  \item Treina $T$ árvores de regressão (CART) com bootstrap
  \item Cada árvore usa subconjunto aleatório de features
  \item Predição final: \textbf{média} das árvores
\end{itemize}
\[
\hat{y}(\mathbf{x}) = \frac{1}{T}\sum_{t=1}^T h_t(\mathbf{x})
\]
\end{block}
\vspace{0.4em}
\begin{exampleblock}{OOB Error}
\small $\approx 36,8\%$ dos exemplos não são sorteados em cada bootstrap --- usados como \textbf{validação gratuita} (Out-of-Bag Error), sem necessidade de conjunto de validação separado.
\end{exampleblock}
\column{0.47\textwidth}
\begin{block}{Vantagens}
\begin{itemize}\small
  \item Robusto a outliers
  \item Pouco ajuste de hiperparâmetros
  \item Importância de features nativa
  \item Paralelo por natureza
\end{itemize}
\end{block}
\vspace{0.4em}
\begin{alertblock}{\small Consulte o Colab}
\small Módulo 4, Notebook com implementação, importância de features e comparação de métricas.
\end{alertblock}
\end{columns}
\end{frame}

\begin{frame}{XGBoost / LightGBM para Regressão}
\begin{columns}
\column{0.42\textwidth}
\begin{block}{Ideia do Boosting}
\small Constrói árvores \textbf{sequencialmente}: cada nova árvore aprende com os erros (resíduos) da anterior.
\[
F_m(\mathbf{x}) = F_{m-1}(\mathbf{x}) + \nu \cdot h_m(\mathbf{x})
\]
\small $\nu$ = taxa de aprendizado (shrinkage)
\end{block}
\vspace{0.4em}
\begin{exampleblock}{\small Vantagem prática}
\small Em dados tabulares com $n < 100k$, XGBoost/LightGBM tipicamente superam redes neurais.
\end{exampleblock}
\column{0.55\textwidth}
\begin{table}
\centering\scriptsize
\begin{tabular}{>{\ttfamily\small}p{3.2cm}p{3.5cm}}
\toprule
\normalfont\textbf{Hiperparâmetro} & \normalfont\textbf{Efeito} \\
\midrule
n\_estimators & Nº de árvores (+ early stop) \\
learning\_rate & $\nu$: menor $\Rightarrow$ mais robusto \\
max\_depth & Profundidade (3--6 ideal) \\
subsample & Fração de exemplos/árvore \\
colsample\_bytree & Fração de features/árvore \\
reg\_alpha (L1) & Regularização nos pesos \\
reg\_lambda (L2) & Regularização nos pesos \\
\bottomrule
\end{tabular}
\end{table}
\begin{alertblock}{\small Colab, Módulo 4}
\small Tuning com Optuna: \texttt{XGBRegressor} e \texttt{LGBMRegressor} com early stopping.
\end{alertblock}
\end{columns}
\end{frame}

% ============================================================
\section{Diagnóstico de Modelos}
% ============================================================

\begin{frame}[t]{Gráficos Diagnósticos}
\centering
\resizebox{0.88\textwidth}{!}{%
\begin{tikzpicture}
% ── TOP-LEFT: Resíduos vs. Ajustados ─────────────────────────
\begin{scope}[xshift=-2cm, yshift=7.8cm]
  \draw[->,minegray!80,thick] (0,-1.7) -- (5.2,-1.7)
    node[right,font=\tiny]{Val.\ Ajustados};
  \draw[->,minegray!80,thick] (0,-1.7) -- (0,1.9)
    node[above,font=\tiny]{Resíduos};
  \draw[mineblue!50, dashed, thick] (0,0) -- (5.2,0);
  \foreach \x/\y in {0.3/0.7,0.7/-1.0,1.1/1.1,1.5/-0.8,2.0/0.5,
                     2.5/-0.7,3.0/1.0,3.5/-0.5,4.0/0.8,4.5/-0.6}
    \filldraw[minegray!70] (\x,\y) circle(0.08);
  \node[font=\small\bfseries,color=mineblue] at (2.6,2.0)
    {Resíduos vs.\ Ajustados};
  \node[font=\tiny,color=minegreen!70!black] at (2.6,-2.0)
    {$\checkmark$ Ideal: dispersão aleatória em torno de zero};
\end{scope}
% ── TOP-RIGHT: Q-Q Plot ──────────────────────────────────────
\begin{scope}[xshift=6.8cm, yshift=7.8cm]
  \draw[->,minegray!80,thick] (0,-1.7) -- (5.2,-1.7)
    node[right,font=\tiny]{Quantis Teóricos};
  \draw[->,minegray!80,thick] (0,-1.7) -- (0,1.9)
    node[above,font=\tiny]{Quantis Amostrais};
  \draw[mineorange, thick] (0.3,-1.5) -- (4.9,1.7);
  \foreach \x/\y in {0.4/-1.45,0.9/-0.95,1.4/-0.50,1.9/-0.05,
                     2.4/0.40,2.9/0.85,3.5/1.20,4.0/1.50,4.6/1.65}
    \filldraw[mineblue!80] (\x,\y) circle(0.08);
  \node[font=\small\bfseries,color=mineblue] at (2.6,2.0) {Q-Q Plot};
  \node[font=\tiny,color=minegreen!70!black] at (2.6,-2.0)
    {$\checkmark$ Ideal: pontos alinhados sobre a reta};
\end{scope}
% ── BOTTOM-LEFT: Scale-Location ──────────────────────────────
\begin{scope}[xshift=-2cm, yshift=0cm]
  \draw[->,minegray!80,thick] (0,0) -- (5.2,0)
    node[right,font=\tiny]{Val.\ Ajustados};
  \draw[->,minegray!80,thick] (0,0) -- (0,3.5)
    node[above,font=\tiny]{$\sqrt{|e_i|}$};
  \draw[minegreen!70, dashed, thick] (0,1.8) -- (5.2,1.8);
  \foreach \x/\y in {0.3/1.7,0.8/1.9,1.3/1.75,1.8/1.8,2.3/1.85,
                     2.9/1.75,3.4/1.8,3.9/1.85,4.5/1.8}
    \filldraw[minegray!70] (\x,\y) circle(0.08);
  \node[font=\small\bfseries,color=mineblue] at (2.6,3.4) {Scale-Location};
  \node[font=\tiny,color=minegreen!70!black] at (2.6,-0.45)
    {$\checkmark$ Ideal: linha aproximadamente horizontal};
\end{scope}
% ── BOTTOM-RIGHT: Distância de Cook ──────────────────────────
\begin{scope}[xshift=6.8cm, yshift=0cm]
  \draw[->,minegray!80,thick] (0,0) -- (5.2,0)
    node[right,font=\tiny]{Índice};
  \draw[->,minegray!80,thick] (0,0) -- (0,3.5)
    node[above,font=\tiny]{Dist.\ de Cook};
  \draw[minered!80, dashed, thick] (0,2.5) -- (5.2,2.5);
  \node[font=\tiny,minered,right] at (5.2,2.5) {limiar};
  \foreach \x/\y in {0.4/0.20,0.9/0.35,1.4/0.18,1.9/0.28,
                     2.4/0.60,2.9/0.22,3.4/0.32,4.0/0.26}
    \draw[mineblue!70,line width=1.5pt] (\x,0) -- (\x,\y);
  \draw[minered, line width=2.5pt] (4.7,0) -- (4.7,2.9);
  \node[font=\tiny\bfseries,minered,above] at (4.7,2.9) {influente!};
  \node[font=\small\bfseries,color=mineblue] at (2.6,3.4) {Distância de Cook};
  \node[font=\tiny] at (2.6,-0.45) {Detecta observações influentes};
\end{scope}
\end{tikzpicture}%
}
\end{frame}

\begin{frame}{O que Procurar em Cada Gráfico}
\begin{columns}
\column{0.48\textwidth}
\begin{block}{Resíduos vs.\ Ajustados}
\begin{itemize}\small
  \item \textcolor{minegreen}{$\checkmark$} Pontos dispersos aleatoriamente
  \item \textcolor{minered}{$\times$} Padrão em funil $\Rightarrow$ heterocedasticidade
  \item \textcolor{minered}{$\times$} Padrão curvo $\Rightarrow$ não-linearidade não capturada
\end{itemize}
\end{block}
\vspace{0.3em}
\begin{block}{Q-Q Plot}
\begin{itemize}\small
  \item \textcolor{minegreen}{$\checkmark$} Pontos sobre a reta $\Rightarrow$ normalidade
  \item \textcolor{minered}{$\times$} Desvios nas caudas $\Rightarrow$ assimetria / caudas pesadas
\end{itemize}
\end{block}
\column{0.48\textwidth}
\begin{block}{Scale-Location}
\begin{itemize}\small
  \item \textcolor{minegreen}{$\checkmark$} Linha horizontal $\Rightarrow$ variância constante
  \item \textcolor{minered}{$\times$} Linha crescente $\Rightarrow$ heterocedasticidade
  \item Solução: transformar $y$ (ex.: $\log y$)
\end{itemize}
\end{block}
\vspace{0.3em}
\begin{block}{Distância de Cook}
\begin{itemize}\small
  \item \textcolor{minegreen}{$\checkmark$} Todos abaixo do limiar
  \item \textcolor{minered}{$\times$} Pontos acima $\Rightarrow$ influentes (investigar)
  \item Limiar usual: $4/n$
\end{itemize}
\end{block}
\end{columns}
\end{frame}

% ============================================================
% REFERÊNCIAS
% ============================================================
\begin{frame}{Referências Bibliográficas}
\begin{thebibliography}{99}
\small
\bibitem{tibshirani1996} Tibshirani, R. (1996). Regression shrinkage and selection via the LASSO. \textit{JRSS-B}, 58(1), 267--288.
\bibitem{breiman2001} Breiman, L. (2001). Random forests. \textit{Machine Learning}, 45(1), 5--32.
\bibitem{chen2016} Chen, T.; Guestrin, C. (2016). XGBoost: A scalable tree boosting system. \textit{KDD 2016}.
\bibitem{friedman2001} Friedman, J. H. (2001). Greedy function approximation: A gradient boosting machine. \textit{Annals of Statistics}, 29(5).
\bibitem{james2021} James, G. et al. (2021). \textit{An Introduction to Statistical Learning} (2nd ed.). Springer.
\end{thebibliography}
\end{frame}

\begin{frame}{}
\centering
\vspace{1em}
{\Large \textcolor{mineblue}{\textbf{Consulte o Material Didático}}}

\vspace{0.8em}
{\large Módulo 4: Técnicas de Regressão}

\vspace{1em}
\begin{block}{\centering Conteúdo Complementar}
\centering
Notebook Google Colab com exemplos práticos de\\
Regressão Linear, Ridge, LASSO, ElasticNet,\\
Random Forest, XGBoost e Diagnóstico de Modelos
\end{block}

\vspace{0.8em}
\small \faGithub \ \href{https://github.com/eltonsarmanho}{github.com/eltonsarmanho}
\end{frame}

\end{document}
